\documentclass[11pt]{article}

\usepackage[utf8]{inputenc}
\usepackage[T1]{fontenc}
\usepackage{lmodern}
\usepackage{amsmath,amssymb,amsthm,mathtools}
\usepackage{stmaryrd}
\usepackage[margin=1in]{geometry}
\usepackage{enumitem}
\usepackage{xcolor}
\usepackage{hyperref}
\usepackage[nameinlink,capitalize]{cleveref}

\hypersetup{
  colorlinks=true,
  linkcolor=blue!60!black,
  citecolor=blue!60!black,
  urlcolor=blue!60!black
}

% theorem environments
\newtheorem{theorem}{Theorem}
\newtheorem{lemma}[theorem]{Lemma}
\newtheorem{proposition}[theorem]{Proposition}
\newtheorem{corollary}[theorem]{Corollary}
\newtheorem{question}[theorem]{Question}
\newtheorem{conjecture}[theorem]{Conjecture}
\theoremstyle{definition}
\newtheorem{definition}[theorem]{Definition}
\newtheorem{example}[theorem]{Example}
\newtheorem{remark}[theorem]{Remark}

% macros
\newcommand{\R}{\mathbb{R}}
\newcommand{\N}{\mathbb{N}}
\newcommand{\E}{\mathbb{E}}
\newcommand{\Prob}{\mathbb{P}}
\newcommand{\A}{\mathcal{A}}
\newcommand{\F}{\mathcal{F}}
\newcommand{\K}{\mathcal{K}}
\newcommand{\T}{\mathcal{T}}
\newcommand{\Homp}{\operatorname{Hom}^{+}}
\newcommand{\Homalg}{\operatorname{Hom}}
\newcommand{\Csem}{\mathcal{C}_{\mathrm{sem}}}
\newcommand{\Ext}{\operatorname{Ext}}
\newcommand{\supp}{\operatorname{supp}}
\newcommand{\ind}{\operatorname{ind}}
\newcommand{\brak}[1]{\left\langle #1\right\rangle}
\newcommand{\avg}[2]{\left\llbracket #1 \right\rrbracket_{#2}}
\newcommand{\q}{\mathfrak q}
\newcommand{\up}{\pi}
\newcommand{\eps}{\varepsilon}

\title{A Blow-Up Generalisation of Forbidden-Subgraph Reasoning in Flag Algebras}
\author{}
\date{\today}

\begin{document}
\maketitle

\begin{abstract}
We prove a general form of the equivalence between two ways of reasoning with a
forbidden-subgraph constraint in Razborov's flag algebras.  The first way passes
to the quotient theory in which the forbidden configurations are axiomatized
away.  The second stays in the ambient graph theory, restricts the unlabelled
root homomorphism to be supported on the forbidden-free class, and then asks for
almost-sure non-negativity under the canonical random extension to a labelled
type.  The equivalence is proved for every hereditary graph class that is closed
under independent blow-ups.  This includes the triangle-free theory, and more
generally the $K_r$-free theories.  The key new point, compared with the
unlabelled case, is a planted blow-up construction: a labelled limit homomorphism
for the quotient theory is embedded with positive probability into the random
extension of a carefully weighted unlabelled blow-up limit.
\end{abstract}

\section{Purpose and statement}

This note addresses the following problem.  Let $\T=\T_{\mathrm{Graph}}$ be the
universal theory of finite simple undirected graphs, and let $\K$ be a hereditary
class of finite graphs.  Equivalently, $\K$ is the class of finite models of the
extension $\T_\K$ obtained from $\T$ by forbidding a set of induced subgraphs.
For a type $\sigma$ that is realised in $\K$, Razborov's interpretation
formalism gives a quotient homomorphism
\[
  \q_\sigma : \A^\sigma[\T] \twoheadrightarrow \A^\sigma[\T_\K].
\]
This is the standard passage from the ambient graph algebra to the algebra in
which the forbidden configurations have been set to zero.  Razborov explicitly
identifies extensions of a theory by forbidding induced subgraphs as quotient
algebras, and he uses $\Homp(\A^\sigma,\R)$ and random extensions of
homomorphisms as the semantic objects for flag algebras \cite[Sections 2.3.3,
3.1--3.2]{Razborov2007}.

The question is whether the quotient statement
\[
  \psi(\q_\sigma f)\ge 0
  \qquad\text{for every }\psi\in\Homp(\A^\sigma[\T_\K],\R)
\]
is equivalent to the following ambient statement: for every unlabelled limit
homomorphism $\phi_0\in\Homp(\A^0[\T],\R)$ that is supported on $\K$, the
canonical random extension $\phi^\sigma$ of $\phi_0$ to type $\sigma$ satisfies
$\phi^\sigma(f)\ge0$ almost surely.  The equivalence is clear when
$\sigma=0$, but not when the labels have positive semantic content.  The main
result below gives a useful general answer.

\begin{definition}[Independent blow-up]
Let $G$ be a finite graph and let $m_v\ge1$ be an integer for each
$v\in V(G)$.  The independent blow-up $G^{\mathbf m}$ is obtained by replacing
$v$ by an independent set $C_v$ of size $m_v$, and by putting all edges between
$C_v$ and $C_w$ whenever $vw\in E(G)$.
\end{definition}

\begin{definition}[Clone-closed hereditary class]
A hereditary graph class $\K$ is \emph{clone-closed} if every independent
blow-up of every $G\in\K$ again belongs to $\K$.
\end{definition}

\begin{example}
The class of $K_r$-free graphs is clone-closed.  Indeed, a clique in an
independent blow-up uses at most one vertex from each clone class, so it
projects to a clique of the same size in the base graph.  In particular, the
triangle-free class is clone-closed.
\end{example}

\begin{remark}[Why this hypothesis is real]
Not every forbidden-induced-subgraph class is clone-closed.  For instance,
$K_2$ is induced-$P_3$-free, but an independent blow-up of $K_2$ with one side
of size two contains an induced $P_3$.  Thus a statement covering arbitrary
hereditary classes needs a different argument or a stronger ambient hypothesis.
The planted blow-up proof below genuinely uses clone-closedness.
\end{remark}

\begin{definition}[$\K$-supported homomorphisms]
Let $\tau$ be a type realised in $\K$.  A homomorphism
$\varphi\in\Homp(\A^\tau[\T],\R)$ is called $\K$-supported if
\[
  \varphi(F)=0
\]
for every $\tau$-flag $F$ whose underlying graph is not in $\K$.
Equivalently, $\varphi$ vanishes on the kernel of the quotient map
$\q_\tau:\A^\tau[\T]\to\A^\tau[\T_\K]$.
\end{definition}

We write
\[
  \brak{\sigma}:=\avg{1_\sigma}{\sigma}\in\A^0[\T]
\]
for the unlabelled density of embeddings of the type $\sigma$.  If
$\phi_0\in\Homp(\A^0[\T],\R)$ and $\phi_0(\brak\sigma)>0$, Razborov's extension
theorem gives a unique probability measure on $\Homp(\A^\sigma[\T],\R)$,
denoted here by
\[
  \phi^\sigma\sim \Ext_\sigma(\phi_0),
\]
characterised by
\begin{equation}\label{eq:extension-expectation}
  \E[\phi^\sigma(g)]
  =
  \frac{\phi_0(\avg{g}{\sigma})}{\phi_0(\brak\sigma)}
  \qquad (g\in\A^\sigma[\T]).
\end{equation}
This is Razborov's ensemble of random homomorphisms rooted at $\phi_0$.

\begin{theorem}[Generalised forbidden-subgraph equivalence]\label{thm:main}
Let $\K$ be a clone-closed hereditary class of finite graphs, and let
$\sigma$ be a non-degenerate type in $\T_\K$.  For every
$f\in\A^\sigma[\T]$, the following are equivalent.
\begin{enumerate}[label=\textup{(\alph*)}, leftmargin=2.8em]
  \item\label{item:quotient}
  For every $\psi\in\Homp(\A^\sigma[\T_\K],\R)$,
  \[
     \psi(\q_\sigma f)\ge0.
  \]
  \item\label{item:ensemble}
  For every $\K$-supported $\phi_0\in\Homp(\A^0[\T],\R)$ with
  $\phi_0(\brak\sigma)>0$, if $\phi^\sigma\sim\Ext_\sigma(\phi_0)$, then
  \[
     \Prob[\phi^\sigma(f)\ge0]=1.
  \]
\end{enumerate}
\end{theorem}

\begin{corollary}[Clique-free theories]\label{cor:cliques}
For every $r\ge3$, every non-degenerate $K_r$-free type $\sigma$, and every
$f\in\A^\sigma[\T_{\mathrm{Graph}}]$, the two conditions in
\Cref{thm:main} are equivalent for $\K=\operatorname{Forb}(K_r)$.  In
particular, the equivalence holds for the triangle-free theory.
\end{corollary}

\section{Preliminaries used in the proof}

We use two standard facts from Razborov's flag algebra formalism.

\begin{theorem}[Razborov representation and extension, used as black boxes]\label{thm:razborov}
Let $\tau$ be a non-degenerate type in a universal graph theory.
\begin{enumerate}[label=\textup{(\roman*)}, leftmargin=2.8em]
  \item Every convergent sequence of $\tau$-flags defines an element of
  $\Homp(\A^\tau,\R)$, and conversely every element of
  $\Homp(\A^\tau,\R)$ is the limit of a convergent sequence of $\tau$-flags.
  \item If $\phi_0\in\Homp(\A^0,\R)$ and $\phi_0(\brak\sigma)>0$, then the
  random extension $\Ext_\sigma(\phi_0)$ exists uniquely and is obtained as
  the weak limit of the finite random-root distributions along any finite
  graph sequence converging to $\phi_0$.
\end{enumerate}
\end{theorem}

Part (i) is Theorem 3.3 of \cite{Razborov2007}; part (ii) is the content of
Theorems 3.5 and 3.12 there.

\begin{lemma}[Quotient homomorphisms are supported homomorphisms]\label{lem:quotient-support}
Composition with $\q_\tau$ identifies $\Homp(\A^\tau[\T_\K],\R)$ with the set
of $\K$-supported homomorphisms in $\Homp(\A^\tau[\T],\R)$.
\end{lemma}

\begin{proof}
The quotient algebra $\A^\tau[\T_\K]$ is obtained from $\A^\tau[\T]$ by setting
to zero exactly those flags whose underlying graph violates the axioms defining
$\K$.  Therefore every homomorphism on the quotient pulls back to a
$\K$-supported homomorphism.  Conversely, any $\K$-supported homomorphism
vanishes on the quotient kernel and hence factors uniquely through the quotient.
\end{proof}

\begin{lemma}[Support passes to random extensions]\label{lem:support-passes}
Let $\phi_0\in\Homp(\A^0[\T],\R)$ be $\K$-supported and suppose
$\phi_0(\brak\sigma)>0$.  If $\phi^\sigma\sim\Ext_\sigma(\phi_0)$, then
$\phi^\sigma$ is $\K$-supported almost surely.
\end{lemma}

\begin{proof}
Let $F$ be a $\sigma$-flag whose underlying graph $G$ is not in $\K$.  The
downward average $\avg{F}{\sigma}$ is a non-negative scalar multiple of the
unlabelled graph $G$.  Since $\phi_0$ is $\K$-supported,
$\phi_0(\avg{F}{\sigma})=0$.  By \eqref{eq:extension-expectation},
\[
  \E[\phi^\sigma(F)]=0.
\]
But $\phi^\sigma(F)\ge0$ pointwise.  Hence $\phi^\sigma(F)=0$ almost surely.
There are only countably many such flags $F$, so intersecting the corresponding
probability-one events shows that $\phi^\sigma$ is $\K$-supported almost surely.
\end{proof}

\section{The planted blow-up lemma}

The hard direction of \Cref{thm:main} is proved by planting the labelled type
inside an unlabelled blow-up with small but positive mass.  The following lemma
is the quantitative form of that idea.

Let $F=(G,\theta)$ be a $\sigma$-flag, with $|\sigma|=k$ and $|G|=n$.  Assume
$n>k$.  Fix a number $0<\rho<1$.  If $k>0$, set
\[
  w_{\theta(i)}=\frac{\rho}{k}\quad(1\le i\le k),
  \qquad
  w_v=\frac{1-\rho}{n-k}\quad(v\notin \operatorname{im}\theta).
\]
If $k=0$, set instead $w_v=1/n$ for every $v\in V(G)$; in this case there are
no planted labelled clusters and the parameter $\rho$ is irrelevant.
For a sufficiently large integer $N$, let $B_N(G,\theta,\rho)$ be an independent
blow-up of $G$ in which the cluster $C_v$ has size approximately $w_vN$.
For this fixed weighted graph define
\[
  \operatorname{err}_N
  :=
  \sum_{v\in V(G)}\left|\frac{|C_v|}{N}-w_v\right|.
\]
By choosing the cluster sizes with $N$ large, this error can be made arbitrarily
small.

\begin{lemma}[Planted blow-up approximation]\label{lem:planted}
Let $m\ge k$ be fixed.  There is a constant $C_m$, depending only on $m$, such
that the following holds.  For every $\sigma$-flag $F_0$ with at most $m$
vertices, every $\sigma$-flag $(G,\theta)$ on $n$ vertices, every sufficiently
large blow-up $B_N=B_N(G,\theta,\rho)$, and every embedding
$\widehat\theta:\sigma\to B_N$ satisfying
$\widehat\theta(i)\in C_{\theta(i)}$ for all $i$, we have
\begin{equation}\label{eq:planted-flag-bound}
  \bigl|p(F_0,(B_N,\widehat\theta))-p(F_0,(G,\theta))\bigr|
  \le
  C_m\left(\rho+\frac1{n-k}+\operatorname{err}_N\right),
\end{equation}
where $\operatorname{err}_N$ is the total rounding error just defined.
Consequently, for every fixed $f\in\A^\sigma[\T]$ there is a constant $C_f$ such
that
\begin{equation}\label{eq:planted-f-bound}
  \bigl|p(f,(B_N,\widehat\theta))-p(f,(G,\theta))\bigr|
  \le
  C_f\left(\rho+\frac1{n-k}+\operatorname{err}_N\right).
\end{equation}
\end{lemma}

\begin{proof}
It is enough to prove \eqref{eq:planted-flag-bound}; the estimate for $f$ then
follows by linearity, since $f$ is a finite linear combination of flags.

Let $\ell=|F_0|\le m$.  To compute $p(F_0,(B_N,\widehat\theta))$, choose
$\ell-k$ additional unlabelled vertices in $B_N$ uniformly, avoiding the labelled
vertices.  Call the sample \emph{good} if:
\begin{enumerate}[label=\textup{(\roman*)}, leftmargin=2.8em]
  \item none of the additional vertices lies in one of the planted labelled
  clusters $C_{\theta(1)},\ldots,C_{\theta(k)}$, and
  \item no two additional vertices lie in the same non-labelled cluster.
\end{enumerate}
The probability that (i) fails is at most $(\ell-k)\rho+O_m(\operatorname{err}_N)$.
Conditional on (i), the non-labelled clusters have equal target weights, and the
rounding error perturbs the induced distribution by at most
$O_m(\operatorname{err}_N)$.  Thus the probability that (ii) fails is at most
$\binom{\ell-k}{2}/(n-k)+O_m(\operatorname{err}_N)$.  Hence the bad probability is bounded by
$C_m(\rho+1/(n-k)+\operatorname{err}_N)$.

On the good event, the projections of the additional vertices to
$V(G)\setminus\operatorname{im}\theta$ are distributed as a uniformly
chosen $(\ell-k)$-subset, up to total variation error $O_m(\operatorname{err}_N)$.  Moreover, the induced
labelled graph in the blow-up is precisely the induced labelled graph obtained
from the projected vertices in $(G,\theta)$, because distinct clone classes have
exactly the adjacencies of their base vertices and the planted labels project to
$\theta(1),\ldots,\theta(k)$.  Thus the two probabilities differ only on the bad
event and by rounding.  This gives \eqref{eq:planted-flag-bound} after enlarging
$C_m$.
\end{proof}

\begin{lemma}[Positive mass of the planted event]\label{lem:positive-mass}
With the notation above, for $N$ sufficiently large the finite random extension
of the unlabelled graph $B_N(G,\theta,\rho)$ to type $\sigma$ assigns probability
at least
\[
  c_{k,\rho}:=
  \begin{cases}
    1, & k=0,\\[2mm]
    \left(\dfrac{\rho}{2k}\right)^k, & k>0,
  \end{cases}
\]
to the event that the random labelled embedding
$\widehat\theta:\sigma\to B_N(G,\theta,\rho)$ satisfies
$\widehat\theta(i)\in C_{\theta(i)}$ for all $i=1,
\ldots,k$.
\end{lemma}

\begin{proof}
For $k=0$ there is nothing to prove.  Assume $k>0$.  Every ordered choice of
vertices $\widehat\theta(i)\in C_{\theta(i)}$ is a model embedding of $\sigma$,
because $\theta$ is a model embedding in $G$ and the blow-up preserves all
relations between distinct base vertices.  For $N$ sufficiently large,
$|C_{\theta(i)}|\ge(\rho/(2k))N$ for each $i$.  Therefore the number of such
planted embeddings is at least $(\rho/(2k))^kN^k$.  The total number of ordered
embeddings of $\sigma$ in $B_N$ is at most $N^k$.  The desired lower bound
follows.
\end{proof}

\section{Proof of the main theorem}

\subsection{Quotient semantics implies ensemble semantics}

Assume \ref{item:quotient}.  Let $\phi_0\in\Homp(\A^0[\T],\R)$ be
$\K$-supported and let $\phi^\sigma\sim\Ext_\sigma(\phi_0)$.  By
\Cref{lem:support-passes}, $\phi^\sigma$ is $\K$-supported almost surely.  By
\Cref{lem:quotient-support}, it therefore descends almost surely to an element
$\overline{\phi}^{\sigma}$ of $\Homp(\A^\sigma[\T_\K],\R)$ satisfying
$\overline{\phi}^{\sigma}(\q_\sigma g)=\phi^\sigma(g)$ for every
$g\in\A^\sigma[\T]$.  Applying \ref{item:quotient} to this descended
homomorphism gives
\[
  \phi^\sigma(f)=\overline{\phi}^{\sigma}(\q_\sigma f)\ge0
\]
almost surely.  Thus \ref{item:ensemble} holds.

\subsection{Ensemble semantics implies quotient semantics}

Assume \ref{item:ensemble}.  We prove \ref{item:quotient} by contradiction.
Suppose that \ref{item:quotient} fails.  Then there exists
$\psi\in\Homp(\A^\sigma[\T_\K],\R)$ with
\[
  \psi(\q_\sigma f)<0.
\]
Pull $\psi$ back along $\q_\sigma$ and still denote the resulting
$\K$-supported homomorphism on $\A^\sigma[\T]$ by $\psi$.  Choose $\alpha>0$ such
that
\begin{equation}\label{eq:negative-margin}
  \psi(f)\le -4\alpha.
\end{equation}

By \Cref{thm:razborov}, applied in the constrained theory $\T_\K$, there is a
convergent sequence of $\sigma$-flags
\[
  F_n=(G_n,\theta_n)\in\F^\sigma[\T_\K],
  \qquad |G_n|\to\infty,
\]
with $p(\cdot,F_n)\to\psi$.  Passing to a tail, we may assume
\begin{equation}\label{eq:finite-negative}
  p(f,F_n)\le -3\alpha
  \qquad\text{for every }n.
\end{equation}
Let $m$ be the maximum number of vertices of a flag appearing in a fixed linear
expansion of $f$.  By \Cref{lem:planted}, choose $0<\rho<1$ so small that
$C_f\rho<\alpha$.  For each $n$, choose a sufficiently large blow-up
\[
  B_n:=B_{N_n}(G_n,\theta_n,\rho)
\]
so that the rounding error in \Cref{lem:planted} is less than $1/n$ and
$|B_n|\to\infty$.  Since $G_n\in\K$ and $\K$ is clone-closed, each $B_n$ also belongs to $\K$.

By compactness, pass to a subsequence, not relabelled, such that the unlabelled
graphs $B_n$ converge in the flag-algebra sense to some
\[
  \phi_0\in\Homp(\A^0[\T],\R).
\]
Since all $B_n$ lie in $\K$, $\phi_0$ is $\K$-supported.  The proof of
\Cref{lem:positive-mass} gives, in particular, a positive lower bound on the
density of embeddings of $\sigma$ in $B_n$ independent of $n$, so
\[
  \phi_0(\brak\sigma)>0.
\]
Thus \ref{item:ensemble} applies to the random extension
$\phi^\sigma\sim\Ext_\sigma(\phi_0)$.

Now consider the finite random extension measure obtained by choosing a random
embedding $\widehat\theta:\sigma\to B_n$.  On the planted event
\[
  \widehat\theta(i)\in C_{\theta_n(i)}\quad(1\le i\le k),
\]
\Cref{lem:planted} and \eqref{eq:finite-negative} give, for all sufficiently
large $n$,
\[
  p(f,(B_n,\widehat\theta))
  \le
  p(f,(G_n,\theta_n))+
  C_f\left(\rho+\frac1{|G_n|-k}+\frac1n\right)
  \le -2\alpha.
\]
By \Cref{lem:positive-mass}, this planted event has probability at least
$c_{k,\rho}>0$.  The finite random extension measures $P_n$ are naturally
measures on the compact product space
$X_\sigma:=[0,1]^{\F^\sigma}$: a sampled embedding
$\widehat\theta$ contributes the point
$p^{(B_n,\widehat\theta)}=(p(F,(B_n,\widehat\theta)))_{F\in\F^\sigma}$.
Thus, with
\[
  C:=\{x\in X_\sigma:x(f)\le -2\alpha\},
\]
we have
\begin{equation}\label{eq:finite-bad-mass}
  P_n(C)\ge c_{k,\rho}
\end{equation}
for all large $n$.  Here $x(f)$ denotes the continuous linear extension of the
coordinate functions to the fixed finite linear combination $f$.

By \Cref{thm:razborov}, the measures $P_n$ converge weakly to the canonical
random extension measure $\Ext_\sigma(\phi_0)$, which is supported on
$\Homp(\A^\sigma[\T],\R)\subseteq X_\sigma$.  The set $C$ is closed in
$X_\sigma$.  The closed-set form of the Portmanteau theorem gives
\[
  \Prob_{\phi^\sigma\sim\Ext_\sigma(\phi_0)}[\phi^\sigma\in C]
  \ge
  \limsup_{n\to\infty} P_n(C)
  \ge c_{k,\rho}>0.
\]
Therefore
\[
  \Prob[\phi^\sigma(f)<0]>0,
\]
contradicting \ref{item:ensemble}.  This contradiction proves
\ref{item:quotient}, and hence \Cref{thm:main}.

\section{Consequences and limitations}

\begin{corollary}[Triangle-free formulation]
Let $\K$ be the class of triangle-free graphs.  For every triangle-free type
$\sigma$ and every $f\in\A^\sigma[\T_{\mathrm{Graph}}]$, the following are
equivalent:
\begin{enumerate}[label=\textup{(\alph*)}, leftmargin=2.8em]
  \item $\psi(\q_\sigma f)\ge0$ for every
  $\psi\in\Homp(\A^\sigma[\T_{K_3\text{-free}}],\R)$;
  \item for every $\phi_0\in\Homp(\A^0[\T_{\mathrm{Graph}}],\R)$ with
  $\phi_0(K_3)=0$ and $\phi_0(\brak\sigma)>0$, the random extension
  $\phi^\sigma\sim\Ext_\sigma(\phi_0)$ satisfies $\phi^\sigma(f)\ge0$ almost
  surely.
\end{enumerate}
\end{corollary}

\begin{proof}
The class of triangle-free graphs is clone-closed by the clique projection
argument in \Cref{cor:cliques}.  Moreover, if $\phi_0(K_3)=0$, then
$\phi_0$ is supported on triangle-free graphs: every graph containing a triangle
has zero $\phi_0$-value by the chain rule and non-negativity of flag densities.
Apply \Cref{thm:main}.
\end{proof}

\begin{remark}[What has been generalised]
The proof does not use any special feature of triangles except the fact that
triangle-free graphs are closed under independent blow-ups.  The same argument
therefore applies to $K_r$-free graphs and to any hereditary graph class with
this clone-closure property.  The proof also handles arbitrary non-degenerate
labelled types $\sigma$, not only the unlabelled algebra.
\end{remark}

\begin{remark}[Why arbitrary hereditary classes remain open for this method]
If $\K$ is not clone-closed, the planted blow-up may leave $\K$.  Then the
unlabelled root homomorphism built in the proof need not satisfy the forbidden
constraints, and the contradiction to ensemble semantics no longer follows.
This does not prove that the equivalence fails for such classes; it only shows
that the present blow-up proof cannot establish it without additional ideas.
A natural next target is to identify weaker replacement hypotheses that still
allow one to realise a labelled quotient homomorphism inside a positive-mass
random extension of an ambient unlabelled homomorphism.
\end{remark}

\begin{thebibliography}{9}

\bibitem{Razborov2007}
A.~A.~Razborov.
\newblock Flag algebras.
\newblock \emph{Journal of Symbolic Logic}, 72(4):1239--1282, 2007.


\end{thebibliography}

\end{document}
